% Section 5.4, from lectures/L05/appendix/c_exercises.md.

\section{Exercises}
\label{c5:sec:exercises}\label{c5:app:c}
These exercises reinforce the concepts from \secref{c5:app:a} and \secref{c5:app:b}.

\begin{aside}
\textbf{How to check your work.} Write each exercise set in its own directory under
\file{lectures/L05/exercises}, with the file names given there, and run \code{make test} from the
repository root. A set’s tests switch on as soon as its files exist, and a set you have not started
is reported as \code{SKIP}. Sets 1 to 3 are one driver growing: start each by copying the previous
set’s directory. The tests of Sets 1 to 3 are ztest applications in
\file{lectures/L05/exercises/test} that build your module against the course’s emulator; Set 4 is
checked by its output, and Set 5 by running your own tests.

The solutions are published in the companion repository, under
\file{lectures/L05/appendix/solutions}, a few days after the lecture. Try each exercise yourself
before you read them.
\end{aside}

Remember \code{source env.sh} in every new terminal.

\subsection{The part: Lumo}
Lumo is a made-up I2C ambient light sensor. It measures continuously: its DATA registers always hold
the latest measurement, and there is nothing to start. Its datasheet, in full:

\begin{booktable}{@{}lllL@{}}
  \thead{Register} & \thead{Name} & \thead{Access} & \thead{Contents} \\
  \midrule
  \code{0x00} & \code{CONFIG} & R/W & Bits 1:0, \code{GAIN}: \code{0} is a gain of 1, \code{1} of 2, \code{2} of 4, \code{3} of 8. Bits 7:2 are reserved, read as 0 and ignore writes. Reset value \code{0x00}. \\
  \code{0x01} & \code{DATA_LSB} & R & The measurement, low byte. \\
  \code{0x02} & \code{DATA_MSB} & R & The measurement, high byte. \\
  \code{0x0A} & \code{PART_ID} & R & Always \code{0x4C} on a genuine part. \\
\end{booktable}

\begin{itemize}
  \item The measurement is an \textbf{unsigned 16-bit count, little-endian}: \code{DATA_LSB} then
    \code{DATA_MSB}.
  \item At a gain of 1, \textbf{one count is 0.25 lux}. At a gain of G, one count is 0.25 / G lux.
  \item A count of \code{0xFFFF} means the part is \textbf{saturated}: the light is at least that bright.
  \item A read of several bytes continues at the following registers, so both DATA bytes can be read in one
    transfer. Writes to read-only registers are ignored.
\end{itemize}

The course provides the part itself: an emulator, in \file{lectures/L05/exercises/test/lumo_emul},
which is a Zephyr module of its own. It plays a Lumo chip for every enabled \code{"qacademy,lumo"}
node on \code{native_sim}’s I2C bus. The tests add it to their builds; to build with it by hand, add
it on the command line:

\begin{shell}
-DZEPHYR_EXTRA_MODULES=$PWD/lectures/L05/exercises/test/lumo_emul
\end{shell}

\noindent Tests control it through that module's own header, \code{<qacademy/lumo_emul.h>}
(\code{lumo_emul_set_part_id()}, \code{lumo_emul_set_counts()} and \code{lumo_emul_get_config()}),
and through the generic sensor emulator backend, with \code{SENSOR_CHAN_LIGHT} in lux.

% ---------------------------------------------------------------------
\exerciseset{A module and a device}
\label{c5:set:1}
Create a directory named \code{exercise1} to work in for this exercise set. The directory is your
module.

\exercise{The module and its Kconfig option}{Code}
\label{c5:ex:1.1}
In this exercise you will create a module that the build can include.

\task{Tasks}
Create these files in \code{exercise1}:

\begin{itemize}
  \item \code{zephyr/module.yml}: a module named \code{qacademy-lumo}, whose CMake and Kconfig are in the
    module’s top directory, and whose \code{dts_root} is the top directory too.
  \item \code{Kconfig}: an option \code{LUMO}, a \code{bool} with the prompt \code{Lumo I2C ambient light
    sensor}, which:
\begin{itemize}
  \item is \code{y} by default,
  \item depends on \code{SENSOR}, and on the devicetree having an enabled \code{"qacademy,lumo"} node,
  \item selects \code{I2C}.
\end{itemize}
  \item \code{CMakeLists.txt}: when \code{CONFIG_LUMO} is enabled, a Zephyr library holding
    \code{src/lumo.c}.
\end{itemize}

\exercise{The binding}{Code}
\label{c5:ex:1.2}

\task{Tasks}
\begin{itemize}
  \item Create \code{dts/bindings/qacademy,lumo.yaml}: a binding for the compatible \code{"qacademy,lumo"},
    which includes \code{sensor-device.yaml} and \code{i2c-device.yaml}, and has no properties of its
    own yet.
  \item Create \code{dts/bindings/vendor-prefixes.txt}, registering the vendor prefix \code{qacademy}.
\end{itemize}

\exercise{The driver’s init}{Code}
\label{c5:ex:1.3}

\task{Tasks}
\begin{itemize}
  \item Create \code{src/lumo_reg.h} with macros for the register addresses, the \code{GAIN} mask, and the
    value of \code{PART_ID}, from the datasheet above.
  \item Create \code{src/lumo.c}, the driver for the compatible \code{"qacademy,lumo"}:
\begin{itemize}
  \item A config struct holding the part’s \code{struct i2c_dt_spec}.
  \item A sensor API table, defined with \code{DEVICE_API(sensor, ...)}, whose \code{sample_fetch} and
    \code{channel_get} functions only return \code{-ENOTSUP} for now.
  \item An init function that:
\begin{itemize}
  \item returns \code{-ENODEV} if the I2C bus is not ready,
  \item reads \code{PART_ID}, and returns the error if the read fails,
  \item returns \code{-ENODEV} if \code{PART_ID} is not \code{0x4C},
  \item and otherwise returns 0.
\end{itemize}
  \item One device per enabled instance, defined with \code{SENSOR_DEVICE_DT_INST_DEFINE()} at
    \code{POST_KERNEL} and \code{CONFIG_SENSOR_INIT_PRIORITY}. The data pointer may be \code{NULL} for
    now.
\end{itemize}
\end{itemize}

The tests in \file{lectures/L05/exercises/test/device} put one part at address \code{0x29}, and a
second at \code{0x2a} whose initialization is deferred, so that a test can give it the wrong
\code{PART_ID} first. To build and run them by hand:

\begin{shell}
west build -p -b native_sim/native/64 -d build/device \
    lectures/L05/exercises/test/device -- \
    -DEXERCISE_DIR=$PWD/lectures/L05/exercises/exercise1
./build/device/zephyr/zephyr.exe
\end{shell}

\noindent The four tests must pass: the option turns itself on, the part is ready, its device has a
sensor API, and the wrong \code{PART_ID} makes \code{device_init()} return \code{-ENODEV}.

\task{Reflection}
\begin{itemize}
  \item Remove \code{vendor-prefixes.txt} and build the tests again by hand. What does the build print? Then
    build once more with \code{-DEXTRA_GEN_EDT_ARGS=--edtlib-Werror} added, as Twister does. What
    happens?
  \item Run \code{west build -d build/device -t initlevels}. Where do your two \code{lumo_init} calls come
    in the list, and what comes just before them? Why must it be in that order?
  \item The deferred part is in the \code{initlevels} list too. Why is it still not ready after boot?
  \item Find \code{DT_HAS_QACADEMY_LUMO_ENABLED} in the build directory. Which file defines it, and why
    could it not exist before you wrote the binding?
\end{itemize}

% ---------------------------------------------------------------------
\exerciseset{Reading light}
\label{c5:set:2}
Copy \code{exercise1} to a directory named \code{exercise2} to work in for this exercise set.

\exercise{Fetch and get}{Code}
\label{c5:ex:2.1}
In this exercise you will implement the sensor API, at the part’s reset gain of 1.

\task{Tasks}
In \code{src/lumo.c}:

\begin{itemize}
  \item Add a data struct holding the count of the last fetch, and pass it to the device.
  \item \code{sample_fetch} shall:
\begin{itemize}
  \item return \code{-ENOTSUP} for any channel other than \code{SENSOR_CHAN_ALL} and
    \code{SENSOR_CHAN_LIGHT},
  \item read both DATA bytes \textbf{in one transfer}, and return the error if the read fails,
  \item store the count, and return 0.
\end{itemize}
  \item \code{channel_get} shall:
\begin{itemize}
  \item return \code{-ENOTSUP} for any channel other than \code{SENSOR_CHAN_LIGHT},
  \item convert the stored count to lux, at 0.25 lux per count, into the \code{struct sensor_value}.
\end{itemize}
  \item Do no I/O in \code{channel_get}.
\end{itemize}

The tests in \file{lectures/L05/exercises/test/conversion} check, among others, that:

\begin{narrowtable}{@{}llll@{}}
  \thead{Count} & \thead{Lux} & \thead{\code{val1}} & \thead{\code{val2}} \\
  \midrule
  \code{0} & 0 & \code{0} & \code{0} \\
  \code{1} & 0.25 & \code{0} & \code{250000} \\
  \code{0x0102} & 64.5 & \code{64} & \code{500000} \\
  \code{0xFFFF} & 16383.75 & \code{16383} & \code{750000} \\
\end{narrowtable}

The tests of Set 1 must still pass.

\textbf{Hint:} \code{sys_get_le16()}, from \code{<zephyr/sys/byteorder.h>}, turns two little-endian
bytes into a \code{uint16_t}. \code{sensor_value_from_micro()} fills in a \code{struct sensor_value}
from millionths.

\task{Reflection}
\begin{itemize}
  \item Replace the single transfer with two \code{i2c_reg_read_byte_dt()} calls, one per DATA byte. Do the
    tests still pass? What could go wrong on a real part that the emulator cannot show?
  \item Compute the lux in \code{int} instead of \code{int64_t}, by removing the cast. Which test fails?
    Above how many counts does the multiplication overflow, and how bright is that?
  \item \code{test_channel_get_returns_the_last_fetch} changes the part’s count between a fetch and a get.
    Why is the old value the right answer?
\end{itemize}

% ---------------------------------------------------------------------
\exerciseset{The gain}
\label{c5:set:3}
Copy \code{exercise2} to a directory named \code{exercise3} to work in for this exercise set.

\exercise{A devicetree property}{Code}
\label{c5:ex:3.1}
In this exercise you will make the gain a property of the devicetree node, written to the part at
init.

\task{Tasks}
\begin{itemize}
  \item In the binding, add a property \code{gain}: an \code{int}, allowed to be only 1, 2, 4 or 8, with a
    default of 1.
  \item In \code{src/lumo.c}:
\begin{itemize}
  \item Add the gain and its \code{GAIN} bits to the config struct. Take the bits from the property’s index
    in the binding’s list of allowed values, at compile time.
  \item In the init function, after checking \code{PART_ID}, write the \code{GAIN} bits to \code{CONFIG},
    without changing the register’s other bits, and return the error if that fails.
  \item In \code{channel_get}, divide by the gain: at a gain of G, one count is 0.25 / G lux.
\end{itemize}
\end{itemize}

The tests in \file{lectures/L05/exercises/test/gain} put three parts on the bus, at gains of 4, 8
and the default, and check what each one’s \code{CONFIG} register holds after boot and what its
counts convert to: one count at a gain of 4 is 0.0625 lux, and \code{0xFFFF} at a gain of 8 is
2047.96875 lux. The tests of Sets 1 and 2 must still pass.

\textbf{Hint:} \code{DT_INST_ENUM_IDX(inst, prop)} is the index of a property’s value in its
binding’s \code{enum} list. \code{i2c_reg_update_byte_dt()} changes only the bits of a register that
a mask selects.

\task{Reflection}
\begin{itemize}
  \item Add \code{gain = <3>;} to a part in an overlay and build the gain tests. What happens, and at which
    step of the build?
  \item Why take the \code{GAIN} bits from \code{DT_INST_ENUM_IDX()}, rather than computing a base-2
    logarithm of the gain at run time?
  \item At a gain of 8, what is the smallest step between two readings, and the brightest light the part can
    measure? What are they at a gain of 1?
\end{itemize}

% ---------------------------------------------------------------------
\exerciseset{An application}
\label{c5:set:4}
Create an application named \code{exercise4} to work in for this exercise set.

\exercise{Reading once a second}{Code}
\label{c5:ex:4.1}

\task{Tasks}
\begin{itemize}
  \item \code{CMakeLists.txt}: add the module in \code{../exercise3} to the build, before finding Zephyr. Do
    \textbf{not} add the emulator.
  \item \code{prj.conf}: enable the sensor subsystem, emulators, and logging.
  \item \code{app.overlay}: a Lumo part labelled \code{light}, at address \code{0x29} on \code{&i2c0}, at the default gain.
  \item \code{src/main.c}, with a log module named \code{light_app} at level \code{INF}:
\begin{itemize}
  \item If the part is not ready, log \code{<name> is not ready} as an error, with the device’s name, and
    return.
  \item Log \code{Reading <name>}.
  \item Then, forever: fetch and get a light reading, log it as \code{<lux> lux} with two decimals, or
    \code{Reading failed: <error>} as an error, and sleep for one second.
\end{itemize}
\end{itemize}

Build it with the emulator, and with the emulator’s sunrise, which changes the light half a second
before each whole second:

\begin{shell}
west build -p -b native_sim/native/64 lectures/L05/exercises/exercise4 -- \
    -DZEPHYR_EXTRA_MODULES=$PWD/lectures/L05/exercises/test/lumo_emul \
    -DCONFIG_EMUL_LUMO_DAYLIGHT=y
./build/zephyr/zephyr.exe -stop_at=6.5
\end{shell}

\noindent The output must contain, in this order, with the emulator’s own log lines in between:

\begin{console}
[00:00:00.000,000] <inf> light_app: Reading lumo@29
[00:00:00.000,000] <inf> light_app: 0.00 lux
[00:00:01.010,000] <inf> light_app: 0.25 lux
[00:00:02.020,000] <inf> light_app: 42.00 lux
[00:00:03.030,000] <inf> light_app: 320.00 lux
[00:00:04.040,000] <inf> light_app: 2750.00 lux
[00:00:05.050,000] <inf> light_app: 12000.00 lux
[00:00:06.060,000] <inf> light_app: 16383.75 lux
\end{console}

\task{Reflection}
\begin{itemize}
  \item Build the application without \code{-DZEPHYR_EXTRA_MODULES=...}. Does it build? What does it print,
    and why?
  \item Build it with the emulator, but with \code{-DCONFIG_EMUL=n} added. What happens, and which
    devicetree node is the error about?
  \item The emulator’s last light level is 30000 lux, but the reading is 16383.75. Why? Set \code{gain =
    <8>} in \code{app.overlay} and run it again: what do the readings become, and what does that say
    about choosing a gain?
\end{itemize}

% ---------------------------------------------------------------------
\exerciseset{Tests of your own}
\label{c5:set:5}
Create a directory named \code{exercise5} to work in for this exercise set: a ztest application that
tests the driver of Exercise Set~\ref{c5:set:3}.

\exercise{A test suite and a Twister scenario}{Code}
\label{c5:ex:5.1}

\task{Tasks}
\begin{itemize}
  \item \code{CMakeLists.txt}: add the module in \code{../exercise3} to the build, and \code{src/main.c} to
    the application. As in Set 4, do not add the emulator.
  \item \code{prj.conf}: enable ztest, the sensor subsystem and emulators.
  \item \code{app.overlay}: two Lumo parts, one at the default gain and one at a gain of 2.
  \item \code{src/main.c}: a suite named \code{lumo}, whose tests check at least:
\begin{itemize}
  \item that both parts are ready,
  \item counts of 0, 1 and \code{0xFFFF} at the default gain,
  \item a count at a gain of 2,
  \item that the gain of 2 was written to the part’s \code{CONFIG} register,
  \item and that an unsupported channel returns \code{-ENOTSUP} from both fetch and get.
\end{itemize}
  \item \code{testcase.yaml}: one scenario, named \code{qacademy.l05.exercise5.lumo}, allowed only on
    \code{native_sim/native/64}, judged by the ztest harness.
\end{itemize}

Build and run it with the emulator:

\begin{shell}
west build -p -b native_sim/native/64 lectures/L05/exercises/exercise5 -- \
    -DZEPHYR_EXTRA_MODULES=$PWD/lectures/L05/exercises/test/lumo_emul
./build/zephyr/zephyr.exe
\end{shell}

\noindent The suite must pass: the output must contain \code{SUITE PASS - 100.00\% [lumo]} and
\code{PROJECT EXECUTION SUCCESSFUL}.

\task{Reflection}
\begin{itemize}
  \item Run the scenario with Twister:
\begin{shell}
west twister -T lectures/L05/exercises/exercise5 -p native_sim/native/64 \
    -x=ZEPHYR_EXTRA_MODULES=$PWD/lectures/L05/exercises/test/lumo_emul
\end{shell}
Where does Twister put the output of your test program? What does it report besides pass and fail?
  \item Break your Set 3 driver on purpose: make \code{channel_get} ignore the gain. Which of your tests
    fail? If none do, which test is missing?
  \item Your tests and the course’s tests both pass. What would you still check on the first board with a
    real Lumo part?
\end{itemize}
